
\chapter{M002 --- Primitive Surface Construction}

\section{Stage 3 --- Implementation Review (Phase 2)}

\subsection{Reviewed Transformation}

Transformation of the conventional Miller covector into its primitive
representation.

\subsection{Mathematical Input}

\begin{itemize}
\item Conventional Miller covector.
\item Verified conventional-to-primitive basis transformation.
\end{itemize}

\subsection{Mathematical Output}

A primitive Miller covector expressed in the primitive reciprocal basis.

\subsection{Algorithmic Role}

This phase transfers the crystallographic orientation into the primitive
reciprocal lattice. It provides the surface orientation used by all
subsequent integer-lattice constructions.

\subsection{Classification}

\begin{itemize}
\item Representation transformation between reciprocal coordinate systems.
\item Preserves crystallographic orientation.
\item Does not create a new scientific surface.
\end{itemize}

\subsection{Objects Preserved}

\begin{itemize}
\item Surface orientation.
\item Miller plane.
\item Crystallographic identity.
\end{itemize}

\subsection{Open Questions for Stage 4}

\begin{enumerate}
\item Is the reciprocal-space transformation mathematically exact?
\item Are orientation invariants preserved for all admissible basis
transformations?
\item Are normalization and sign conventions explicitly documented?
\end{enumerate}
